\documentclass[11pt,a4paper]{article}
\usepackage[utf8]{inputenc}
\usepackage[margin=1in]{geometry}
\usepackage{amsmath,amssymb}
\usepackage{listings}
\usepackage{xcolor}
\usepackage{hyperref}
\usepackage{graphicx}
\usepackage{fancyhdr}
\usepackage{titlesec}

% Colors
\definecolor{codegreen}{rgb}{0,0.6,0}
\definecolor{codegray}{rgb}{0.5,0.5,0.5}
\definecolor{codepurple}{rgb}{0.58,0,0.82}
\definecolor{backcolour}{rgb}{0.95,0.95,0.92}

% Code listing style
\lstdefinestyle{mikoshi}{
    backgroundcolor=\color{backcolour},   
    commentstyle=\color{codegreen},
    keywordstyle=\color{magenta},
    numberstyle=\tiny\color{codegray},
    stringstyle=\color{codepurple},
    basicstyle=\ttfamily\footnotesize,
    breakatwhitespace=false,         
    breaklines=true,                 
    captionpos=b,                    
    keepspaces=true,                 
    numbers=left,                    
    numbersep=5pt,                  
    showspaces=false,                
    showstringspaces=false,
    showtabs=false,                  
    tabsize=2
}
\lstset{style=mikoshi}

% Headers
\pagestyle{fancy}
\fancyhf{}
\rhead{MikoshiLang Reference Manual}
\lhead{Mikoshi Ltd}
\rfoot{Page \thepage}

% Title formatting
\titleformat{\section}{\Large\bfseries}{\thesection}{1em}{}
\titleformat{\subsection}{\large\bfseries}{\thesubsection}{1em}{}

\hypersetup{
    colorlinks=true,
    linkcolor=blue,
    filecolor=magenta,      
    urlcolor=cyan,
}

\title{\Huge\textbf{MikoshiLang}\\\Large Reference Manual \& User Guide}
\author{Mikoshi Ltd}
\date{Version 1.1.5 --- \today}

\begin{document}

\maketitle
\thispagestyle{empty}

\vspace{2cm}

\begin{abstract}
MikoshiLang is a symbolic computation language for Python, inspired by Wolfram Language. It provides 1,342 functions across 20+ domains including calculus, linear algebra, number theory, statistics, chemistry, visualization, and more. Unlike traditional computational tools, MikoshiLang combines natural language processing with deterministic code execution — AI proposes, the engine proves. No hallucinated math.
\end{abstract}

\vfill

\begin{center}
\textbf{Try it online:} \url{https://mikoshi.co.uk/mikoshilang/console}\\
\textbf{GitHub:} \url{https://github.com/DarrenEdwards111/MikoshiLang}\\
\textbf{PyPI:} \texttt{pip install mikoshilang}
\end{center}

\newpage

\tableofcontents

\newpage

\section{Introduction}

MikoshiLang is a symbolic computation language designed for Python, providing Wolfram-style syntax with deterministic evaluation. Built on SymPy, it extends Python's mathematical capabilities with pattern matching, rule-based transformations, and natural language processing.

\subsection{Key Features}

\begin{itemize}
    \item \textbf{1,342 functions} across calculus, algebra, number theory, statistics, chemistry, units, visualization
    \item \textbf{Natural language mode} — type plain English, get verified results
    \item \textbf{Pattern matching \& rules} — symbolic manipulation with custom transformations
    \item \textbf{Jupyter kernel} — interactive notebooks with rich output
    \item \textbf{Visualization} — 2D/3D plotting with matplotlib \& plotly
    \item \textbf{Chemistry} — equation balancing, molecular mass, element data
    \item \textbf{Unit conversion} — 50+ physical units with dimensional analysis
    \item \textbf{No hallucinations} — AI translates natural language to code, engine verifies
\end{itemize}

\subsection{Installation}

\begin{lstlisting}[language=bash]
# Standard installation
pip install mikoshilang

# With visualization support
pip install mikoshilang[visualization]

# All optional dependencies
pip install mikoshilang[all]
\end{lstlisting}

\subsection{Quick Start}

\begin{lstlisting}[language=Python]
from mikoshilang import parse_and_eval

# Basic arithmetic
parse_and_eval("2 + 2")
# => 4

# Symbolic calculus
parse_and_eval("D[x^3 + 2*x, x]")
# => 3*x^2 + 2

# Equation solving
parse_and_eval("Solve[x^2 - 4 == 0, x]")
# => {-2, 2}

# Plotting
parse_and_eval("Plot[sin(x), {x, 0, 6.28}]")
# => <base64 PNG image>
\end{lstlisting}

\section{Syntax}

MikoshiLang uses Wolfram-inspired syntax with square brackets for function calls and curly braces for lists.

\subsection{Basic Expressions}

\begin{lstlisting}
# Arithmetic
2 + 3 * 4       # => 14
2^10            # => 1024
Sqrt[16]        # => 4

# Variables
x + 2*y - z     # => x + 2*y - z

# Lists
{1, 2, 3, 4}    # => [1, 2, 3, 4]
Range[1, 10]    # => [1, 2, 3, ..., 10]
\end{lstlisting}

\subsection{Function Calls}

\begin{lstlisting}
# Single argument
Sin[Pi/2]       # => 1
Exp[1]          # => E

# Multiple arguments
GCD[48, 18]     # => 6
Mod[17, 5]      # => 2

# Nested calls
Sqrt[Abs[-16]]  # => 4
\end{lstlisting}

\subsection{Pattern Matching}

\begin{lstlisting}
# Define a rule
rule = x^2 + 2*x + 1 -> (x + 1)^2

# Apply rule
ReplaceAll[x^2 + 2*x + 1, rule]
# => (x + 1)^2
\end{lstlisting}

\section{Core Functions}

\subsection{Calculus}

\subsubsection{Differentiation}

\begin{lstlisting}
D[x^3, x]                    # First derivative
D[x^3 + 2*x^2, x, 2]         # Second derivative
\end{lstlisting}

\textbf{Available:} \texttt{D}, \texttt{Derivative}, \texttt{PartialD}, \texttt{Gradient}, \texttt{Divergence}, \texttt{Curl}, \texttt{Laplacian}

\subsubsection{Integration}

\begin{lstlisting}
Integrate[x^2, x]            # Indefinite integral
Integrate[x^2, {x, 0, 1}]    # Definite integral (0 to 1)
\end{lstlisting}

\textbf{Available:} \texttt{Integrate}, \texttt{NIntegrate}, \texttt{LineIntegral}, \texttt{DoubleIntegral}, \texttt{TripleIntegral}

\subsubsection{Limits}

\begin{lstlisting}
Limit[sin(x)/x, x, 0]        # => 1
Limit[1/x, x, Infinity]      # => 0
\end{lstlisting}

\textbf{Available:} \texttt{Limit}, \texttt{Series}, \texttt{TaylorSeries}, \texttt{LaurentSeries}

\subsection{Algebra}

\subsubsection{Solving Equations}

\begin{lstlisting}
Solve[x^2 - 4 == 0, x]       # => {-2, 2}
Solve[x^2 + x + 1 == 0, x]   # Complex solutions
NSolve[cos(x) == x, x]       # Numerical solver
\end{lstlisting}

\textbf{Available:} \texttt{Solve}, \texttt{NSolve}, \texttt{DSolve}, \texttt{RSolve}, \texttt{Roots}, \texttt{FindRoot}

\subsubsection{Simplification}

\begin{lstlisting}
Simplify[(x^2 - 1)/(x - 1)]  # => x + 1
Expand[(x + 1)^3]            # => x^3 + 3*x^2 + 3*x + 1
Factor[x^2 - 4]              # => (x - 2)*(x + 2)
\end{lstlisting}

\textbf{Available:} \texttt{Simplify}, \texttt{FullSimplify}, \texttt{Expand}, \texttt{Factor}, \texttt{Apart}, \texttt{Together}, \texttt{Cancel}, \texttt{Collect}

\subsection{Linear Algebra}

\begin{lstlisting}
# Determinant
Det[{{1, 2}, {3, 4}}]        # => -2

# Eigenvalues
Eigenvalues[{{1, 2}, {2, 1}}]  # => [-1, 3]

# Matrix multiplication
Dot[{{1, 2}, {3, 4}}, {{5}, {6}}]  # => {{17}, {39}}
\end{lstlisting}

\textbf{Available:} \texttt{Det}, \texttt{Inverse}, \texttt{Transpose}, \texttt{Eigenvalues}, \texttt{Eigenvectors}, \texttt{Dot}, \texttt{Cross}, \texttt{MatrixPower}, \texttt{MatrixRank}, \texttt{Trace}, \texttt{Norm}

\subsection{Number Theory}

\begin{lstlisting}
Prime[100]                   # 100th prime number
PrimeQ[17]                   # True (17 is prime)
FactorInteger[60]            # => {{2, 2}, {3, 1}, {5, 1}}
GCD[48, 18]                  # => 6
LCM[12, 18]                  # => 36
\end{lstlisting}

\textbf{Available:} \texttt{Prime}, \texttt{PrimeQ}, \texttt{NextPrime}, \texttt{PrimePi}, \texttt{FactorInteger}, \texttt{Divisors}, \texttt{GCD}, \texttt{LCM}, \texttt{EulerPhi}, \texttt{MoebiusMu}, \texttt{Fibonacci}, \texttt{Lucas}

\subsection{Statistics}

\begin{lstlisting}
Mean[{1, 2, 3, 4, 5}]        # => 3
Median[{1, 2, 3, 4, 5}]      # => 3
StandardDeviation[{1, 2, 3}] # => sqrt(2/3)
Variance[{1, 2, 3, 4}]       # => 5/4
\end{lstlisting}

\textbf{Available:} \texttt{Mean}, \texttt{Median}, \texttt{Mode}, \texttt{Variance}, \texttt{StandardDeviation}, \texttt{Quantile}, \texttt{Quartiles}, \texttt{Correlation}, \texttt{Covariance}, \texttt{Skewness}, \texttt{Kurtosis}

\section{Visualization}

MikoshiLang supports 2D and 3D plotting with matplotlib (static) or plotly (interactive).

\subsection{2D Plotting}

\begin{lstlisting}
# Simple plot
Plot[sin(x), {x, 0, 6.28}]

# Multiple functions
Plot[{sin(x), cos(x)}, {x, 0, 6.28}]

# Parametric plot
ParametricPlot[{cos(t), sin(t)}, {t, 0, 6.28}]

# Polar plot
PolarPlot[1 + cos(theta), {theta, 0, 6.28}]
\end{lstlisting}

\subsection{3D Plotting}

\begin{lstlisting}
# 3D surface
Plot3D[sin(x)*cos(y), {x, -3, 3}, {y, -3, 3}]

# Contour plot
ContourPlot[x^2 + y^2, {x, -2, 2}, {y, -2, 2}]

# Vector field
VectorFieldPlot[{-y, x}, {x, -2, 2}, {y, -2, 2}]
\end{lstlisting}

\subsection{Interactive Plots}

\begin{lstlisting}
# Use plotly for interactive 3D
Plot3D[x^2 + y^2, {x, -2, 2}, {y, -2, 2}, interactive=True]
\end{lstlisting}

\section{Chemistry}

\subsection{Equation Balancing}

\begin{lstlisting}
BalanceEquation["H2 + O2 -> H2O"]
# => 2H2 + O2 -> 2H2O

BalanceEquation["C6H12O6 -> CO2 + H2O"]
# => C6H12O6 -> 6CO2 + 6H2O
\end{lstlisting}

\subsection{Molecular Mass}

\begin{lstlisting}
MolecularMass["H2O"]         # => 18.015 g/mol
MolecularMass["C6H12O6"]     # => 180.156 g/mol (glucose)
MolecularMass["NaCl"]        # => 58.443 g/mol
\end{lstlisting}

\subsection{Element Data}

\begin{lstlisting}
Element["Fe"]
# => {
#   name: Iron,
#   number: 26,
#   mass: 55.845,
#   electronConfiguration: [Ar] 3d6 4s2,
#   category: transition metal
# }
\end{lstlisting}

\section{Unit Conversion}

MikoshiLang supports 50+ physical units across length, mass, time, temperature, energy, power, pressure, and more.

\begin{lstlisting}
Convert[100, "miles", "km"]      # => 160.934 km
Convert[32, "fahrenheit", "celsius"]  # => 0 C
Convert[1, "atm", "pascal"]      # => 101325 Pa
Convert[1, "kWh", "joules"]      # => 3.6e6 J
\end{lstlisting}

\textbf{Available units:} meters, km, miles, feet, inches, kg, pounds, ounces, seconds, hours, days, celsius, fahrenheit, kelvin, joules, calories, kWh, watts, horsepower, atm, bar, pascal, liters, gallons

\section{Natural Language Mode}

MikoshiLang includes an AI-powered natural language interface. Type plain English, get verified results.

\subsection{How It Works}

\begin{enumerate}
    \item User types natural language query
    \item AI translates to MikoshiLang code
    \item Code is executed by the deterministic engine
    \item Result is verified and returned
\end{enumerate}

\subsection{Examples}

\begin{lstlisting}
# Natural language input
"What is the derivative of x cubed?"
# AI translates to: D[x^3, x]
# Engine returns: 3*x^2

"Solve x squared equals 4"
# AI translates to: Solve[x^2 == 4, x]
# Engine returns: {-2, 2}

"What is the molecular mass of water?"
# AI translates to: MolecularMass["H2O"]
# Engine returns: 18.015 g/mol
\end{lstlisting}

\subsection{Web Console}

Try the interactive web console at:

\url{https://mikoshi.co.uk/mikoshilang/console}

Features auto-detection — if your input looks like natural language, it routes to the NL endpoint automatically.

\section{API Usage}

\subsection{Python Library}

\begin{lstlisting}[language=Python]
from mikoshilang import parse_and_eval, evaluate
from mikoshilang.parser import parse

# Parse and evaluate
result = parse_and_eval("D[x^2, x]")
print(result)  # => 2*x

# Parse only (returns AST)
ast = parse("x^2 + 2*x + 1")
print(ast)  # => BinOp(...)

# Evaluate parsed AST
result = evaluate(ast)
print(result)  # => x^2 + 2*x + 1
\end{lstlisting}

\subsection{REST API}

MikoshiLang is deployed at \texttt{mikoshi.co.uk} with a REST API:

\begin{lstlisting}[language=bash]
# Evaluate expression
curl -X POST https://mikoshi.co.uk/api/mikoshilang/eval \
  -H "Content-Type: application/json" \
  -d '{"input": "D[x^3, x]", "mode": "parse"}'

# Natural language
curl -X POST https://mikoshi.co.uk/api/mikoshilang/nl \
  -H "Content-Type: application/json" \
  -d '{"input": "What is the integral of x squared?"}'

# Chemistry
curl -X POST https://mikoshi.co.uk/api/mikoshilang/chemistry \
  -H "Content-Type: application/json" \
  -d '{"action": "balance", "input": "H2 + O2 -> H2O"}'
\end{lstlisting}

\subsection{Jupyter Kernel}

\begin{lstlisting}[language=bash]
# Install kernel
python -m mikoshilang.kernel install

# Start Jupyter
jupyter notebook

# Select "MikoshiLang" kernel
\end{lstlisting}

\section{Function Reference}

MikoshiLang provides 1,342 functions across 20+ domains. Below is a categorized listing.

\subsection{Calculus (67 functions)}

\texttt{D, Derivative, PartialD, Grad, Gradient, Div, Divergence, Curl, Laplacian, Integrate, NIntegrate, Limit, Series, TaylorSeries, LaurentSeries, FourierSeries, FourierTransform, InverseFourierTransform, LaplaceTransform, InverseLaplaceTransform, ZTransform, InverseZTransform, DSolve, NDSolve, Residue, LineIntegral, SurfaceIntegral, VolumeIntegral, DoubleIntegral, TripleIntegral, ...}

\subsection{Algebra (89 functions)}

\texttt{Solve, NSolve, DSolve, RSolve, Roots, FindRoot, Simplify, FullSimplify, Expand, Factor, Apart, Together, Cancel, Collect, Coefficient, CoefficientList, PolynomialQ, Degree, Exponent, ReplaceAll, ReplaceRepeated, Rule, Pattern, Blank, BlankSequence, ...}

\subsection{Linear Algebra (76 functions)}

\texttt{Det, Determinant, Inverse, Transpose, Eigenvalues, Eigenvectors, CharacteristicPolynomial, MinimalPolynomial, Dot, Cross, Norm, Normalize, MatrixPower, MatrixExp, MatrixLog, MatrixRank, Trace, DiagonalMatrix, IdentityMatrix, ZeroMatrix, RandomMatrix, HilbertMatrix, VandermondeMatrix, JordanForm, SchurDecomposition, QRDecomposition, LUDecomposition, CholeskyDecomposition, SVD, PseudoInverse, ...}

\subsection{Number Theory (94 functions)}

\texttt{Prime, PrimeQ, IsPrime, NextPrime, PreviousPrime, PrimePi, PrimeFactors, FactorInteger, Divisors, DivisorSigma, EulerPhi, MoebiusMu, LiouvilleLambda, Fibonacci, Lucas, Binomial, Multinomial, Factorial, DoubleFactorial, Subfactorial, GCD, LCM, ExtendedGCD, ChineseRemainder, Mod, PowerMod, JacobiSymbol, LegendreSymbol, KroneckerSymbol, ...}

\subsection{Statistics (112 functions)}

\texttt{Mean, Median, Mode, Variance, StandardDeviation, Quantile, Quartiles, InterquartileRange, Range, Correlation, Covariance, Skewness, Kurtosis, GeometricMean, HarmonicMean, TrimmedMean, CentralMoment, Moment, RandomVariate, PDF, CDF, InverseCDF, NormalDistribution, UniformDistribution, ExponentialDistribution, PoissonDistribution, BinomialDistribution, ChiSquareDistribution, StudentTDistribution, FDistribution, ...}

\subsection{Trigonometry (48 functions)}

\texttt{Sin, Cos, Tan, Cot, Sec, Csc, ArcSin, ArcCos, ArcTan, ArcCot, ArcSec, ArcCsc, Sinh, Cosh, Tanh, Coth, Sech, Csch, ArcSinh, ArcCosh, ArcTanh, ArcCoth, ArcSech, ArcCsch, Haversine, InverseHaversine, TrigExpand, TrigReduce, TrigSimplify, TrigToExp, ExpToTrig, ...}

\subsection{Special Functions (97 functions)}

\texttt{Gamma, LogGamma, PolyGamma, DiGamma, TriGamma, Beta, Zeta, HurwitzZeta, Erf, Erfc, Erfi, BesselJ, BesselY, BesselI, BesselK, AiryAi, AiryBi, HermiteH, LaguerreL, LegendreP, ChebyshevT, ChebyshevU, JacobiP, GegenbauerC, SphericalHarmonicY, EllipticK, EllipticE, EllipticF, EllipticPi, ...}

\subsection{Discrete Math (68 functions)}

\texttt{Permutations, Combinations, Multinomial, Subsets, Tuples, IntegerPartitions, SetPartitions, StirlingS1, StirlingS2, BellNumber, CatalanNumber, BernoulliB, EulerE, HarmonicNumber, GeneralizedHarmonicNumber, PartitionsP, PartitionsQ, ...}

\subsection{Logic \& Sets (34 functions)}

\texttt{And, Or, Not, Xor, Nand, Nor, Implies, Equivalent, ForAll, Exists, Union, Intersection, Complement, SymmetricDifference, Subset, Element, Cardinality, PowerSet, CartesianProduct, ...}

\subsection{String Operations (52 functions)}

\texttt{StringLength, StringReverse, StringJoin, StringSplit, StringReplace, StringTake, StringDrop, ToUpperCase, ToLowerCase, StringContainsQ, StringStartsQ, StringEndsQ, StringCount, StringPosition, Characters, FromCharacterCode, ToCharacterCode, ...}

\subsection{List Manipulation (89 functions)}

\texttt{Length, First, Last, Rest, Most, Part, Take, Drop, Append, Prepend, Insert, Delete, Sort, Reverse, Flatten, Partition, Riffle, Join, Union, Intersection, Complement, Select, Map, Apply, Fold, FoldList, Nest, NestList, Table, Range, ...}

\subsection{Graph Theory (64 functions)}

\texttt{Graph, DirectedGraph, WeightedGraph, AdjacencyMatrix, IncidenceMatrix, EdgeList, VertexList, VertexCount, EdgeCount, VertexDegree, InDegree, OutDegree, PathGraph, CycleGraph, CompleteGraph, StarGraph, WheelGraph, GridGraph, TreeGraph, ShortestPath, GraphDistance, GraphDiameter, GraphRadius, ConnectedComponents, StronglyConnectedComponents, WeaklyConnectedComponents, ArticulationPoints, Bridges, ChromaticNumber, ChromaticPolynomial, MaximumFlow, MinimumCut, SpanningTree, MinimumSpanningTree, EulerianPath, HamiltonianPath, ...}

\subsection{Chemistry (118 elements + functions)}

\texttt{Element, MolecularMass, BalanceEquation, AtomicNumber, AtomicMass, ElectronConfiguration, Electronegativity, MeltingPoint, BoilingPoint, Density, H, He, Li, Be, B, C, N, O, F, Ne, Na, Mg, Al, Si, P, S, Cl, Ar, K, Ca, Fe, Cu, Zn, Ag, Au, U, ...}

\subsection{Units (50+ units)}

\texttt{Convert, Meter, Kilometer, Mile, Foot, Inch, Kilogram, Pound, Ounce, Second, Hour, Day, Celsius, Fahrenheit, Kelvin, Joule, Calorie, KilowattHour, Watt, Horsepower, Atmosphere, Bar, Pascal, Liter, Gallon, ...}

\subsection{Visualization (9 functions)}

\texttt{Plot, Plot3D, ParametricPlot, PolarPlot, ContourPlot, VectorFieldPlot, Heatmap, AnimatePlot, Interactive3D}

\subsection{Numerical Methods (43 functions)}

\texttt{NIntegrate, NSolve, FindRoot, FindMinimum, FindMaximum, NMinimize, NMaximize, NDSolve, InterpolatingFunction, Interpolation, LinearInterpolation, CubicSplineInterpolation, BSplineInterpolation, Fit, LinearFit, PolynomialFit, NonlinearModelFit, ...}

\subsection{Optimization (37 functions)}

\texttt{Minimize, Maximize, NMinimize, NMaximize, LinearProgramming, ConvexOptimization, ConstrainedMin, ConstrainedMax, GradientDescent, NewtonMethod, GoldenSectionSearch, ...}

\subsection{Control Theory (28 functions)}

\texttt{TransferFunction, StateSpaceModel, Poles, Zeros, SystemGain, BodePlot, NyquistPlot, RootLocus, StabilityMargin, Controllability, Observability, LQR, PIDController, ...}

\subsection{Signal Processing (46 functions)}

\texttt{FFT, InverseFFT, DFT, IDFT, DCT, IDCT, Spectrogram, Periodogram, Correlate, Convolve, ButterworthFilter, ChebyshevFilter, BesselFilter, EllipticFilter, ...}

\subsection{Time \& Date (24 functions)}

\texttt{Now, Today, DateString, DateList, DayName, MonthName, DayOfWeek, DayOfYear, LeapYearQ, DateDifference, DatePlus, TimeZone, UnixTime, ...}

\subsection{Random \& Probability (38 functions)}

\texttt{Random, RandomInteger, RandomReal, RandomChoice, RandomSample, SeedRandom, RandomVariate, RandomPrime, RandomGraph, RandomMatrix, ...}

\subsection{Financial Math (21 functions)}

\texttt{PresentValue, FutureValue, NPV, IRR, Annuity, TimeValue, CompoundInterest, ContinuouslyCompounded, BondPrice, YieldToMaturity, ...}

\section{Advanced Topics}

\subsection{Custom Rules}

Define your own transformation rules:

\begin{lstlisting}[language=Python]
from mikoshilang import parse_and_eval

# Define a custom simplification rule
rule = "sin(x)^2 + cos(x)^2 -> 1"

# Apply it
expr = "sin(x)^2 + cos(x)^2 + 2*x"
result = parse_and_eval(f"ReplaceAll[{expr}, {rule}]")
# => 1 + 2*x
\end{lstlisting}

\subsection{Pattern Matching}

\begin{lstlisting}
# Match any power
Pattern[_, x^_]

# Match specific forms
Pattern[a_ + b_, ...]
\end{lstlisting}

\subsection{Performance Tips}

\begin{itemize}
    \item Use \texttt{NSolve} for numerical roots instead of symbolic \texttt{Solve}
    \item Cache repeated computations
    \item Use \texttt{N[expr]} to convert symbolic to numeric
    \item Simplify expressions before integration/differentiation
    \item For large matrices, use sparse representations
\end{itemize}

\section{Troubleshooting}

\subsection{Common Errors}

\textbf{SyntaxError: Invalid expression}

Ensure square brackets for function calls and proper operator syntax.

\textbf{NameError: Function not found}

Check spelling and capitalization. MikoshiLang functions are case-sensitive (\texttt{Sin} not \texttt{sin}).

\textbf{ValueError: Cannot parse expression}

Verify parentheses/brackets are balanced.

\subsection{Getting Help}

\begin{itemize}
    \item Web console: \url{https://mikoshi.co.uk/mikoshilang/console}
    \item GitHub Issues: \url{https://github.com/DarrenEdwards111/MikoshiLang/issues}
    \item Documentation: \url{https://github.com/DarrenEdwards111/MikoshiLang/blob/main/README.md}
\end{itemize}

\section{Contributing}

MikoshiLang is open-source (Apache 2.0). Contributions welcome:

\begin{enumerate}
    \item Fork the repository
    \item Create a feature branch
    \item Add functions to \texttt{mikoshilang/extended*.py}
    \item Write tests in \texttt{tests/}
    \item Submit a pull request
\end{enumerate}

\section{License}

MikoshiLang is licensed under the Apache License 2.0.

\begin{verbatim}
Copyright 2025 Mikoshi Ltd

Licensed under the Apache License, Version 2.0 (the "License");
you may not use this file except in compliance with the License.
You may obtain a copy of the License at

    http://www.apache.org/licenses/LICENSE-2.0
\end{verbatim}

\section{Acknowledgments}

MikoshiLang is built on:
\begin{itemize}
    \item \textbf{SymPy} — symbolic mathematics
    \item \textbf{matplotlib} — 2D plotting
    \item \textbf{plotly} — interactive visualization
    \item \textbf{numpy} — numerical arrays
\end{itemize}

Inspired by Wolfram Language and designed for the Python ecosystem.

\vfill

\begin{center}
\Large\textbf{Built by Mikoshi Ltd}\\
\normalsize
\url{https://mikoshi.co.uk}
\end{center}

\end{document}
